
\section{2. Related Work}

\subsection{2.1 Efficient Attention}

\subsubsection{2.1.1 Sparse Attention}

Longformer (Beltagy et al., 2020) restricts attention to a sliding window 
plus global tokens, achieving $\mathcal{O}(L \sqrt{L})$. BigBird (Zaheer et 
al., 2020) adds random connections. Primal Attention is \textbf{dense} in content 
but \textbf{low-rank} in structure.

\subsubsection{2.1.2 Kernel Attention}

Performer (Choromanski et al., 2021) approximates softmax via random Fourier 
features, achieving $\mathcal{O}(L)$ with a fixed number of random features. 
Primal Attention uses landmarks instead of random features.

\subsubsection{2.1.3 Low-Rank Attention}

Linformer (Wang et al., 2020) projects $K, V$ to a fixed low rank via learned 
projections. Nystr mformer (Xiong et al., 2021) uses Nystr m landmark 
approximation with the same landmarks for $Q$ and $K$. Primal Attention 
\textbf{breaks this symmetry}.

\subsubsection{2.1.4 Performer vs Nystr m}

\begin{table}[H]
\centering
\begin{tabular}{llll}
\toprule
\textbf{Method} & \textbf{Approximation} & \textbf{Symmetric?} & \textbf{Trainable?} \\
\midrule
Performer & Random Fourier features & Yes & Yes \\
Nystr mformer & Symmetric landmarks & Yes & Yes \\
**Primal Attention** & Asymmetric landmarks & **No** & **Yes** \\
\bottomrule
\end{tabular}
\end{table}

\subsection{2.2 Nystr m Methods}

\subsubsection{2.2.1 Classical Nystr m}

Williams \& Seeger (2001) approximate a Gram matrix $K_{L \times L}$ as 
$C W^{-1} C^\top$. This requires $C \in \mathbb{R}^{L \times m}$ and 
$W \in \mathbb{R}^{m \times m}$.

\subsubsection{2.2.2 Nystr mformer}

Xiong et al. (2021) apply Nystr m to attention by selecting $m$ landmarks 
via \textbf{mean pooling} with stride $L/m$, then using the same landmarks for 
both $Q$ and $K$. Primal Attention generalizes this by allowing separate 
landmarks.

\subsubsection{2.2.3 Randomized Nystr m}

Drineas \& Mahoney (2005) select landmarks randomly. Li et al. (2021) use 
k-means clustering. Primal Attention uses \textbf{segment pooling} for simplicity 
and determinism.

\subsection{2.3 Pseudoinverse Computation}

\subsubsection{2.3.1 SVD-based Pseudoinverse}

Standard: $A^+ = V \Sigma^+ U^\top$. Cost: $\mathcal{O}(m^3)$ per layer, 
which is the dominant cost.

\subsubsection{2.3.2 Schulz Iteration}

Schulz (1933) proposed the iterative scheme $X_{t+1} = X_t (2I - A X_t)$, 
which converges quadratically to $A^{-1}$. We adapt this for pseudoinverse 
with a slight modification.

\subsubsection{2.3.3 Newton-Schulz in Deep Learning}

Bernstein \& Newhouse (2024) use Newton-Schulz iteration for orthogonalization 
in optimizers. We use it for attention pseudoinverse.

\subsection{2.4 Position of Our Work}

\begin{table}[H]
\centering
\begin{tabular}{lll}
\toprule
\textbf{Axis} & \textbf{Nystr mformer} & \textbf{Primal Attention} \\
\midrule
Landmark selection & Mean pooling (shared) & Segment pooling (asymmetric) \\
Kernel symmetry & Symmetric & **Asymmetric** \\
Pseudoinverse & SVD-based & **Schulz iteration** \\
Differentiable & Yes & Yes \\
GQA support & No & Yes \\
RoPE support & Partial & Full \\
\bottomrule
\end{tabular}
\end{table}

\textbf{Key distinction.} Primal Attention's asymmetric formulation allows the 
model to use different summaries for $Q$ and $K$, capturing the different 
roles of these projections. This is not possible in standard Nystr m 
approximation.